%------------------------------------------------------------------------------%

\usepackage{integracao_e_series}

\title{Derivadas}

%------------------------------------------------------------------------------%
\begin{document}

\addtitle

%------------------------------------------------------------------------------%
\section{Derivadas}

%------------------------------------------------------------------------------%
\begin{frame}{Derivada de uma Função Real}
  A \emph{Derivada} de uma função real $f(x)$ em relação
  a variável $x$
  \pause
  é a função $f'$\\
  cujo valor em $x$ é
  \[
    f'(x) = \lim_{h\to 0} \frac{\,f(x+h)-f(x)\,}{h}
  \]
  desde que o limite exista.
\end{frame}

%------------------------------------------------------------------------------%
\begin{frame}{Interpretação Geométrica da Derivada}
  A \emph{derivada} de $f$ no ponto $x$ é \\[3mm]
  \pause
  o \emph{coeficiente angular}
  da reta tangente ao gráfico de $f$ no ponto $x$
\end{frame}

%------------------------------------------------------------------------------%
\begin{frame}{Função Derivada}
  A \emph{função derivada} é \\[3mm]
  \pause
  a função \(f'(x)\) que retorna o valor da derivada de $f$ em cada ponto $x$
\end{frame}

%------------------------------------------------------------------------------%
\begin{frame}{Algumas Funções Derivada}
\begin{columns}
\column{0.5\linewidth}
\begin{gather*}
  \onslide<+->{ \frac{dc}{dx} = 0 \\[3mm]}
  \onslide<+->{ \frac{dx}{dx} = 1 \\[3mm]}
  \onslide<+->{ \frac{d}{dx} x^n = n x^{n-1}}
\end{gather*}
\column{0.5\linewidth}
\begin{gather*}
  \onslide<+->{ \frac{d}{dx} e^x = e^x \\[3mm]}
  \onslide<+->{ \frac{d}{dx} \ln\abs{x} = \frac{1}{x}}
\end{gather*}
\end{columns}
\end{frame}

%------------------------------------------------------------------------------%
\begin{frame}{Algumas Funções Derivada}
  \begin{columns}
    \column{0.5\linewidth}
    \begin{gather*}
      \onslide<+->{ \frac{d}{dx} \sin(x) = \cos(x) \\[3mm]}
      \onslide<+->{ \frac{d}{dx} \cos(x) = -\sin(x) \\[3mm]}
      \onslide<+->{ \frac{d}{dx} \tan(x) =  \sec^2(x) \\[3mm]}
    \end{gather*}
    \column{0.5\linewidth}
    \begin{gather*}
      \onslide<+->{ \frac{d}{dx} \cot(x) = -\csc^2(x) \\[3mm]}
      \onslide<+->{ \frac{d}{dx} \sec(x) =  \sec(x) \tan(x) \\[3mm]}
      \onslide<+->{ \frac{d}{dx} \csc(x) = -\csc(x) \cot(x)}
    \end{gather*}
  \end{columns}
\end{frame}

%------------------------------------------------------------------------------%
\begin{frame}{Propriedades da Derivada -- Regas de Derivação}
  $f$ e $g$ funções deriváveis, $c$ e $r$ constantes
  \begin{columns}
    \column{0.5\linewidth}
    \begin{enumerate}[<+->]
      \item \( \big(cf \big)' = cf' \) \\[3mm]
      \item \( \big(f+g\big)' = f' + g' \)\\[3mm]
      \item \( \big(f-g\big)' = f' - g' \)\\[3mm]
      \item \( \big(fg \big)' = f'g + fg' \)
    \end{enumerate}
    \column{0.5\linewidth}
    \vspace{\baselineskip}
    \begin{enumerate}[<+->]
      \setcounter{enumi}{4}
      \item \( \left(\frac{f}{\,g\,}\right)' = \frac{f'g - fg'}{g^2} \)\\[7mm]
      \item Se \(h(x) = f\big(g(x)\big) \) então  \[ h'(x) = f'\big(g(x)\big)\,g'(x) \]
    \end{enumerate}
  \end{columns}
\end{frame}

%------------------------------------------------------------------------------%
\section{Lista Mínima}

%------------------------------------------------------------------------------%
\begin{frame}{Lista Mínima}
Apostila ``\href{https://sites.google.com/view/prof-luis-dafonseca/gradua\%C3\%A7\%C3\%A3o/Integrais-e-Series}{Integração e Séries}''
\begin{enumerate}
  \item Estudar atentamente texto do Capítulo 1
\end{enumerate}
\vfill
Atenção: A prova é baseada no livro, não nas apresentações
\end{frame}

%------------------------------------------------------------------------------%
\end{document}
%------------------------------------------------------------------------------%
